\section{Prefix-complete aggregate reassembly}
\label{sec:aggregate}

Let \(\Theta_X\) be a finite native-key set.  For each
\(\theta\in\Theta_X\), let \(\rho_{\theta,X}\ge0\) be the literal
origin, \(N_{\theta,X}\ge1\) the returned length, and
\[
 b_{\theta,X}(n),\qquad 1\le n\le N_{\theta,X},
\]
the complete local coefficient after retaining the affine M\"obius
factors, periodic factors, smooth weight, residue and coprimality
masks, and the additive determinant phase when present.  Define
\begin{align}
 S_{\theta,X}(N)
 &:=
 \sum_{n\le N}b_{\theta,X}(n),
 \label{eq:fiber-sum}\\
 L_{\theta,X}^{(\rho)}(N)
 &:=
 \sum_{n\le N}
 \frac{b_{\theta,X}(n)}{\rho_{\theta,X}+n}.
\label{eq:fiber-anchored-primitive}
\end{align}
The symbol \(\rho_{\theta,X}\) is an interval origin; it is unrelated
to the Fourier frequency \(r\) used in the TPC-93 low window.

\begin{theorem}[Exact anchored packet identity]
\label{thm:aggregate-identity}
Suppose the literal raw packet has the exact reassembly
\begin{equation}
 Z_X=
 \sum_{\theta\in\Theta_X}
 c_{\theta,X}
 S_{\theta,X}(N_{\theta,X})
 +E_{\mathrm{ret},X}.
\label{eq:literal-reassembly}
\end{equation}
Then
\begin{equation}
 \boxed{
 Z_X=
 \sum_{\theta\in\Theta_X}
 c_{\theta,X}N_{\theta,X}
 \cT_{\rho_{\theta,X},N_{\theta,X}}
 (L_{\theta,X}^{(\rho)})
 +E_{\mathrm{ret},X}.}
\label{eq:aggregate-signed}
\end{equation}
In particular,
\begin{align}
 |Z_X|
 \le{}&
 \sum_{\theta\in\Theta_X}
 |c_{\theta,X}|N_{\theta,X}
 \left|
 \cT_{\rho_{\theta,X},N_{\theta,X}}
 (L_{\theta,X}^{(\rho)})
 \right|
 \notag\\
 &+|E_{\mathrm{ret},X}|.
\label{eq:aggregate-absolute}
\end{align}
\end{theorem}

\begin{proof}
Apply \cref{thm:anchored-identity} to every finite fiber in
\eqref{eq:literal-reassembly}.  No limiting interchange is required.
The inequality follows after the exact outer sum has been formed.
\end{proof}

\begin{remark}[Literal TPC-93 convention]
\label{rem:tpc93-convention}
In a TPC-93 application the native key must include the low Fourier
frequency, polarization, exact-content child, masks, and interval
component.  The sign in the exact relation
\[
 Z_X=-\operatorname{Low}_X+\operatorname{Tail}_X,
\]
together with the reconstruction multiplier and each
frequency-dependent unit phase, stays in \(c_{\theta,X}\) or
\(b_{\theta,X}\) according to its literal location.  The returned
remainder is then
\[
 E_{\mathrm{ret},X}=\operatorname{Tail}_X,
\]
not an invented content error.  The complete source--child inverse
and multiplicity reconstruction established in TPC-93 supplies this
\(\Lone\) attachment \citep{WangTPC93}.
\end{remark}

\begin{corollary}[Quantitative zero-mode certificate]
\label{cor:zero-mode-certificate}
At the critical physical scales, if
\begin{align}
 &\sum_{\theta\in\Theta_X}
 |c_{\theta,X}|N_{\theta,X}
 \left|
 \cT_{\rho_{\theta,X},N_{\theta,X}}
 (L_{\theta,X}^{(\rho)})
 \right|
 +|E_{\mathrm{ret},X}|
 \notag\\
 &\hspace{35mm}\le
 X^{-\eta_Z+o(1)}Q_X^2,
\label{eq:aggregate-target}
\end{align}
then
\[
 |Z_X|\le X^{-\eta_Z+o(1)}Q_X^2.
\]
\end{corollary}

\begin{proof}
Use \eqref{eq:aggregate-absolute}.
\end{proof}

\begin{remark}[Signed aggregate is the minimal target]
\label{rem:signed-minimal}
The absolute certificate \eqref{eq:aggregate-target} is sufficient,
not necessary.  The exact minimal object is the signed outer sum in
\eqref{eq:aggregate-signed}.  Taking absolute values before native
outer reassembly may lose a fixed power.  Any argument exploiting
outer cancellation must still retain every native coefficient and
origin.
\end{remark}

\begin{remark}[Bounded fibers]
\label{rem:bounded-fibers}
A fiber with \(N_{\theta,X}=O(1)\) has no internal asymptotic
parameter.  It remains in \eqref{eq:aggregate-signed} exactly and
must be controlled by its coefficient mass or by cancellation after
outer reassembly.  Deleting it because a long-fiber theorem does not
apply changes \(Z_X\).
\end{remark}

\begin{proposition}[Coefficient-mass promotion rule]
\label{prop:coefficient-mass}
Suppose a subset \(\Theta_X^\star\) satisfies
\[
 \left|
 \cT_{\rho_{\theta,X},N_{\theta,X}}
 (L_{\theta,X}^{(\rho)})
 \right|
 \le\varepsilon_X
 \quad(\theta\in\Theta_X^\star).
\]
Its absolute contribution is at most
\[
 \varepsilon_X
 \sum_{\theta\in\Theta_X^\star}
 |c_{\theta,X}|N_{\theta,X}.
\]
Therefore a uniform \(o(1)\) defect is useful only relative to the
actual weighted outer mass.
\end{proposition}

\begin{proof}
This is the corresponding part of
\eqref{eq:aggregate-absolute}.
\end{proof}
